% Figure: the developing entrance region of pipe flow. The wall boundary layers
% grow from the inlet and merge on the axis at the end of the entrance length
% L_e, after which the profile is the fully developed parabola. The lower panel
% shows the wall shear stress decaying from its high entrance value to the
% constant developed value. Schematic; the curve for wall shear is
% tau_w/tau_dev = 1 + a*exp(-b*x/L_e).
\begin{figure}[htbp]
\centering
\begin{tikzpicture}
% Upper panel: pipe with growing boundary layers.
\begin{scope}
% Pipe walls.
\draw[enggray, thick] (0,1.2) -- (8,1.2);
\draw[enggray, thick] (0,-1.2) -- (8,-1.2);
% Boundary-layer edges growing from each wall, merging near x = 5.
\draw[engblue, very thick]
  (0,1.2) .. controls (2.5,0.55) and (4.0,0.12) .. (5,0);
\draw[engblue, very thick]
  (0,-1.2) .. controls (2.5,-0.55) and (4.0,-0.12) .. (5,0);
% Fully developed region.
\draw[engblue, very thick, dashed] (5,0) -- (8,0);
% Inlet uniform profile arrows.
\foreach \y in {-0.9,-0.45,0,0.45,0.9} {
  \draw[engred, -{Latex}] (0.05,\y) -- (0.75,\y);
}
% Developed parabolic profile at exit.
\draw[engred, very thick]
  (7.2,-1.2) .. controls (7.9,-0.6) and (7.9,0.6) .. (7.2,1.2);
\foreach \y in {-0.9,-0.45,0,0.45,0.9} {
  \draw[engred, -{Latex}] (7.2,\y) -- ({7.2 + 0.7*(1 - (\y/1.2)*(\y/1.2))},\y);
}
% Entrance-length marker.
\draw[enggray, {Latex}-{Latex}] (0,-1.6) -- (5,-1.6);
\node[font=\scriptsize, enggray, anchor=north] at (2.5,-1.6) {entrance length $L_e$};
\node[font=\scriptsize, engblue, anchor=south] at (2.3,0.75) {growing boundary layer};
\node[font=\scriptsize, anchor=south] at (6.5,1.25) {fully developed};
\node[font=\scriptsize, engred, anchor=west] at (0.05,1.5) {uniform inlet};
\end{scope}
\end{tikzpicture}
\caption[The entrance region of pipe flow]{Downstream of an inlet the flow is
nearly uniform across the bore with thin wall boundary layers. Over the
entrance length $L_e$ the layers thicken and merge on the axis, after which the
profile is the fully developed parabola and stops changing with distance. The
wall shear stress is largest at the inlet, where the boundary layer is thin and
its wall gradient steep, and it falls to the constant developed value once the
layers merge, so the pressure drop over the entrance region exceeds the
developed prediction. The laminar entrance length grows with Reynolds number,
$L_e/D\approx0.05\,\mathrm{Re}_D$, while the turbulent one is far shorter and
nearly Reynolds-independent, which is why entrance effects are chiefly a
laminar and short-pipe concern.}
\label{fig:vol03ch12:entrance-region}
\end{figure}
